\section{A radial threshold and automatic bounded density}\label{sec:radial}

We now exclude the unbounded-density alternative for a stable scalar
solution $v:\R^7\to(-1,1)$ with $D>0$. Throughout this section $m=6$.
The inputs are the reach inequality and its positive-foot transfer from
Proposition~\ref{prop:reach-calculus}, the sharp interface estimates from
Section~\ref{sec:interface}, and the fourth-order bound
\eqref{eq:M1}. All are uniform under translation. We will use this
uniformity to choose the center of a large ball after choosing its radius.

The proof has two parts. First we compare reach and stability using the
same surface amplitude. The comparison forces a positive radial integral
to be negligible. We then transfer the energy in the selected ball to
regular zero-surface area. Jointly injective normal tubes show that most
of this area has small reach, where the radial integral has a definite
lower bound. These conclusions are incompatible.

\subsection{The common amplitude and the comparison to be proved}

Recall that $I=I_P$ is the normal average with normalization
$\sigma=4\sqrt2/3$, and $d$ is the penalized reach. Put
$K(R)=1+\sup_xM_R(x)$ and suppose $K(R)\to\infty$. For each large $R$
choose a center with $M_R\ge K(R)/2$, then translate that center to the
origin. Define
\[
 \rho=(R^2+|x|^2)^{1/2},\qquad \tau=\log\rho,\qquad
 h=\rho^{-2},\qquad \phi=h^2=\rho^{-4},\qquad
 L=\frac{R}{\sqrt{K(R)}}.
\]
The length $L$ is large but small compared with $R$. Indeed
$K(R)\le1+CR$ gives $L\ge c_0\sqrt R$, while
\begin{equation}\label{eq:radial-scale-separation}
 \sup_{\rho\ge R}\frac{L+\log\rho}{\rho}
 \le\frac1{\sqrt{K(R)}}+\frac{\log R}{R}\longrightarrow0.
\end{equation}
This separation requires no rate of divergence for $K$.

Fix
\[
 \vartheta=.02,\qquad c=.87,\qquad q_*=2.2,\qquad
 \alpha=.04,\qquad\beta_0=.01.
\]
Choose a smooth nondecreasing $f_*:\R\to(0,1]$ such that
\[
 0\le(\log f_*)'\le q_*,\qquad f_*(s)=1\quad(s\ge1),\qquad
 f_*(s)=C_*e^{q_*s}\quad(s\le-1).
\]
For $t>0$ set
\[
 b(t)=\bigl(t^{1-\vartheta}-L^{1-\vartheta}\bigr)_+^{1/(1-\vartheta)},
 \qquad W_*(\tau,t)=t^{-\vartheta}f_*(c\tau-b(t)),\qquad
 \mathcal F(\tau,t)=\int_t^\infty W_*(\tau,u)\dd u.
\]
Below the reach scale $L$, the weight is simply $t^{-\vartheta}$.
Past the transition $b(t)\simeq c\tau$ it decays exponentially.
We evaluate $W_*$ and $\mathcal F$ at $t=d(z)$ and write
$\lambda=(\log f_*)'(c\tau-b(d))$. The symbols $W_*$ and $f_*$
are distinct from the potential $W$ and the decreasing reach multiplier
$f$ of Section~\ref{sec:growth}.

The amplitude will vanish near zeros where the defect is not yet small.
Define
\[
 Z=\{z\in\Sigma:D(z)\ge\rho(z)^{-\alpha}\}.
\]
Choose a smooth ambient $m_H:\R^7\to[0,1]$, zero when
$\dist(x,Z)<\rho(x)^{\beta_0}$ and one when
$\dist(x,Z)>4\rho(x)^{\beta_0}$, with uniformly bounded first and second
derivatives. One obtains it by convolving an intermediate distance
cutoff on a fixed small scale; the gaps between the displayed constants
preserve these two requirements for large $R$. Set $m_H=1$ if $Z$ is
empty. Every singular zero has $D=1$ and lies in $Z$.
Let $m_B$ be smooth, equal to one on $B_B$, zero outside $B_{2B}$,
with $|\nabla^km_B|\le C_kB^{-k}$ for $k=1,2$, and put $m_c=m_Hm_B$.
All supported surface expressions therefore lie in the regular surface.
Moreover, Lemma~\ref{lem:reach-small-defect} gives
\[
 d_{\min}(R):=\inf_{\{m_H\ne0\}\cap\Sigma}d\longrightarrow\infty
 \quad(\inf\varnothing:=+\infty).
\]
The common amplitudes and mass are
\[
 V=h\sqrt{W_*},\qquad \psi=m_cV,\qquad
 \mathcal L=\int_\Sigma I\psi^2\dd S,\qquad
 T_r=1-|\nabla_\Sigma\rho|^2\ge R^2/\rho^2.
\]
When an amplitude or its derivatives is written on all of $\Sigma$, it
is extended by zero through the collar where $m_H=0$.

The integrated reach inequality, proved in
Lemma~\ref{lem:radial-reach-budget}, will have right side
$.73\mathcal L$. An ambient test interpolating the values of $\psi$
between adjacent sheets will give
\[
 .895\mathcal L\le1.005\int_\Sigma|\nabla_\Sigma\psi|^2\dd S
       +\text{negligible errors}.
\]
The initial surface estimate \eqref{eq:L3} does not give this much
interaction mass when $a^2$ is small compared with $I$. Nor do the
numbers $.73$ and $.895$ alone yield a contradiction. Expanding the
gradient for this amplitude separates the radial
contribution from the reach-gradient contribution. Their coefficients
will be compared after the two inequalities have been proved; the strict
difference leaves a positive multiple of
\[
 \int_\Sigma\phi m_c^2\mathcal F\,\frac{T_r}{\rho^2}\dd S.
\]
This identifies the quantity that the final area argument must bound
from below.

We measure errors against $\mathcal A=K(R)L^{1-\vartheta}$.
An error is \emph{negligible} if its absolute value is at most
$e(R)\mathcal A+e_R(B)$, where $e(R)\to0$ and, for each fixed $R,L$,
$e_R(B)\to0$ as $B\to\infty$. All estimates are uniform in the chosen
center. The limits are always taken in the order $B\to\infty$ at fixed
$R$, followed by $R\to\infty$.

\subsection{Weight identities and fixed analytic choices}

We now fix the analytic tolerances used in the comparison. First choose
the small normal-interpolation allowances quantified in
Lemma~\ref{lem:radial-switch-cost}. Choose the penalty $\Gamma$ so that
$1\le G_1<1.001$ and $\eta_0<10^{-4}$ in
\eqref{eq:C8}--\eqref{eq:C9}, and set $\kappa=\sqrt2/G_1$.
Prescribe the tolerances in \eqref{eq:L1}, \eqref{eq:L12}, and
\eqref{eq:L13}; then fix a sufficiently long normal interval $[-P,P]$.
In particular $|H|\le.355I$, while the two sharp interaction tolerances
may be at most $.001$. The middle-phase estimate holds at its prescribed
tolerance on sufficiently large clean windows. Finally choose the defect
threshold and the other constants in \eqref{eq:M1}. The interpolation
half-width $A_{\rm sw}=3$, depth constant $M$, and reach truncation constant
$M_0$ are fixed before the final increase of $R$. None of these choices
uses a lower bound on the speed at which $K(R)$ diverges.

\Needspace{3\baselineskip}
\begin{lemma}[Weight identities]\label{lem:radial-weight}
Uniformly for $\tau\ge\log R$ and $t\ge t_{\min}$, where $t_{\min}>0$ is
any fixed lower bound for the reach on the good set,
\begin{equation}\label{eq:R1}
\begin{split}
 \mathcal F_\tau
 &=(1+o_R(1))\,cW_*
       \left(\frac{t}{c\tau\vee b(t)}\right)^\vartheta,\\
 \frac{\mathcal F_\tau+\mathcal F}{W_*}&\le C(L+\tau),\qquad
 \frac{\mathcal F}{W_*}\ge\frac1{q_*+\vartheta/t_{\min}},\\
 \mathcal F&\le C(L+\tau)^{1-\vartheta}.
\end{split}
\end{equation}
On the regular zero surface the following is a distributional measure
inequality:
\begin{equation}\label{eq:R2}
\begin{split}
 \Delta_\Sigma\mathcal F\ge{}&
 \mathcal F_\tau\Delta_\Sigma\tau
 -c^2\vartheta J_1|\nabla_\Sigma\tau|^2
 -W_*(\Delta_\Sigma d)_{\rm upper}\\
 &+\frac{\vartheta W_*}{d}|\nabla_\Sigma d|^2
 +\frac{W_*\lambda}{b'}|\nabla_\Sigma(c\tau-b(d))|^2,\\
 J_1={}&\int_{b(d)}^\infty
 u^{-\vartheta-1}f_*'(c\tau-u)\dd u
 \le C\mathcal F_\tau/\tau.
\end{split}
\end{equation}
The last square is defined to be zero where $b'=0$.
\end{lemma}

\begin{proof}
For $t>L$,
\[
 b'=(b/t)^\vartheta,\qquad t^{-\vartheta}\dd t=b^{-\vartheta}\dd b.
\]
Since $c\tau>1$ for large $R$, the part $t\le L$ contributes nothing to
$\mathcal F_\tau$. Thus
\[
 \mathcal F_\tau
 =c\int_{b(t)}^\infty u^{-\vartheta}f_*'(c\tau-u)\dd u.
\]
The total mass of the measure $f_*'(c\tau-u)\dd u$ in this integral is
$f_*(c\tau-b(t))$. If $b(t)\le c\tau+1$, this mass is concentrated
at $u=c\tau+O(1)$ with an exponentially decaying right tail; if
$b(t)>c\tau+1$, it is concentrated at $u=b(t)+O(1)$ with the same
property. Its support is contained in $u\ge c\tau-1$ in the first case.
Comparing $u^{-\vartheta}$ with
$(c\tau\vee b(t))^{-\vartheta}$ proves the first estimate uniformly;
the same comparison with $u^{-\vartheta-1}$ proves the bound for $J_1$.

The logarithmic slope satisfies
\begin{equation}\label{eq:radial-weight-slopes}
 0\le-\partial_t\log W_*=\vartheta/t+\lambda b'\le
 q_*+\vartheta/t_{\min},\qquad 0\le\partial_\tau\log W_*\le cq_*.
\end{equation}
Integration of the resulting exponential lower bound for
$W_*(\tau,t+s)/W_*(\tau,t)$ gives the lower ratio in \eqref{eq:R1}.
The same argument starting at the current $t$ gives
$\mathcal F/W_*\ge1/(q_*+\vartheta/t)$. In particular, on any region
where $d\ge d_{\min}(R)\to\infty$, the ratio is at least
$1/(2.2+o_R(1))$ uniformly.
The relation $t\le C(L+b(t))$, splitting the integrals at $C(L+\tau)$,
and the exponential tail of $f_*$ give the two upper bounds.
In the exponentially decaying region one can also use
\[
 \mathcal F/W_*\le C(t/b(t))^\vartheta\le C
 \quad\text{if }t>C(L+\tau).
\]
This observation will allow polynomial factors in $d$ to be discarded
far beyond the switching region.

For completeness the derivatives responsible for the gradient square are
\[
 \mathcal F_t=-W_*,\quad
 \mathcal F_{tt}=W_*(\vartheta/t+\lambda b'),\quad
 \mathcal F_{\tau t}=-cW_*\lambda,\quad
 \mathcal F_{\tau\tau}
 =c^2\bigl(b^{-\vartheta}f_*'(c\tau-b)-\vartheta J_1\bigr).
\]
Since $b^{-\vartheta}f_*=W_*/b'$, their quadratic expression in
$(\nabla_\Sigma\tau,\nabla_\Sigma d)$ is exactly the last two terms
of \eqref{eq:R2}, minus $c^2\vartheta J_1|\nabla_\Sigma\tau|^2$.
Near $t=L$ we have $f_*'=0$; all terms involving $\lambda/b'$ therefore
vanish on a neighborhood of this apparent singularity.
The reach is locally semiconcave, with nonpositive singular Laplacian
part. Multiplication by $\mathcal F_t=-W_*<0$ changes this part to a
nonnegative measure. The smooth calculation for its almost-everywhere
density consequently gives \eqref{eq:R2}; one can equivalently obtain it
by local concave smoothing and passage to the distributional limit.
\end{proof}

\subsection{Removing rough regions and summing the errors}

The fourth-order estimate pays for the collar of $Z$ and the regions
where the reach inequality has a large curvature error. The next lemma
records the two sums needed for those payments and the separate decay
needed to remove the outer cutoff at fixed $R$.

\begin{lemma}[Summable errors]\label{lem:radial-errors}
Write $s_j=2^jR$, $j\ge0$; a shell $\rho\asymp s_j$ may be enlarged by
any fixed factor. The good surface area and ambient volume within
$C s_j^{\beta_0}$ of $Z$ on this shell are bounded by
$C K(s_j)s_j^{m-4+2\alpha+n\beta_0}$. Each of the following two shell
bounds has negligible sum:
\begin{equation}\label{eq:R3}
 C K(s)s^{-2+2\alpha+n\beta_0}(1+\log s)^C
       (L+\log s)^{1-\vartheta},
\end{equation}
\begin{equation}\label{eq:R4}
 C K(s)e^{C\sqrt{\log(K(s)+1)}}s^{-2}
       (L+\log s)^{3-\vartheta}.
\end{equation}
Moreover, for some $\epsilon'>0$ and every fixed $R,L$,
\begin{equation}\label{eq:R5}
 \int_{\{\rho\asymp s\}\cap(\Sigma\setminus X)}W_*\dd S
 \le C_Rs^{m-\epsilon'}\qquad(s\longrightarrow\infty).
\end{equation}
\end{lemma}

\begin{proof}
Take a uniformly separated net on $Z$ in an enlarged shell.
Outside a unit neighborhood of $X$, each such center has a regular
surface disk of fixed radius on which $D\ge c s^{-\alpha}$, by
\eqref{eq:diffuse-local}. Its contribution to $\int D^2$ is at least
$c s^{-2\alpha}$. Bounded overlap and \eqref{eq:M1} bound the number
of these centers by $C K(s)s^{m-4+2\alpha}$. Centers near $X$ obey
the stronger bad-net bound from the same estimate. Enlargement by
$Cs^{\beta_0}$ costs at most $Cs^{n\beta_0}$, using unit ambient
volume and unit good-area bounds.

For \eqref{eq:R3} use $K(2^jR)\le C2^jK(R)$ and
$-2+2\alpha+n\beta_0=-1.85$. After division by $\mathcal A$, its sum
is bounded by
\[
 C R^{-1.85}(1+\log R)^C
 \sum_{j\ge0}2^{-.85j}(1+j)^C=o(1).
\]
Here $(L+\log R+j)/L\le C(1+j)$ since $L\ge c_0\sqrt R$.
For \eqref{eq:R4}, the elementary inequality
$\sqrt{a+b}\le\sqrt a+\sqrt b$ gives instead
\begin{align*}
 \frac1{\mathcal A}\sum_{j\ge0}
 C K(s_j)e^{C\sqrt{\log(K(s_j)+1)}}s_j^{-2}
       (L+\log s_j)^{3-\vartheta}
 &\le C\frac{L^2}{R^2}
 e^{C\sqrt{\log(K(R)+1)}}
 \sum_{j\ge0}2^{-j}e^{C\sqrt{j+1}}(1+j)^C\\
 &\le C\frac{e^{C\sqrt{\log(K(R)+1)}}}{K(R)}
 \longrightarrow0.
\end{align*}
The last limit uses only $K(R)\to\infty$.

To prove \eqref{eq:R5}, fix a sufficiently small $\epsilon>0$.
At fixed $L$, $b(t)\ge(1-\epsilon)t-C_{\epsilon,L}$ for all $t>0$.
Thus, on the good set,
\[
 W_*\le C_Rs^{cq_*}E^p,\qquad
 p=\frac{q_*(1-\epsilon)}{\kappa}<2.
\]
Hölder's inequality, \eqref{eq:M1}, and the good-area bound give
$\int_{\rho\asymp s}E^p\le C K(s)s^{m-2p}$.
Using $K(s)\le Cs$ proves \eqref{eq:R5} with
\[
 \epsilon'=q_*\left(\frac{2(1-\epsilon)}{\kappa}-c\right)-1>0.
\]
For example the limiting value at $\epsilon=0$, $\kappa=\sqrt2$ is
greater than $.197$, so the fixed choices above leave strict room.
\end{proof}

Here are several consequences used without repeating the summation.
Derivative costs in the rough collar are bounded by \eqref{eq:R3}.
For reach derivatives this follows from \eqref{eq:C7} and the weight:
a factor $e^{C\sqrt{\log(d+1)}}$ is bounded by $C_\xi d^\xi$ for
every fixed $\xi>0$, while $W_*\le d^{-\vartheta}$. More explicitly,
$b'\le1$, $\lambda\le q_*$, $|\nabla\tau|\le\rho^{-1}$, and
$K(d)\le1+Cd$ imply
\[
 W_*|\nabla_\Sigma\log W_*|^2
 \le C d^{-\vartheta}
 \left[\rho^{-2}+(1+d^{-2})e^{C\sqrt{\log(d+1)}}\right]
 \le C.
\]
The last constant is independent of $L$; the negative power of $d$
absorbs the subpolynomial factor, including near $d=L$.
Along logarithmic normal segments, the corresponding metric
and volume factors are also bounded by such powers. The same reasoning
applies to terms involving $\mathcal F$ using \eqref{eq:R1}.

Outer cutoff errors without mean curvature are
$C_R(1+\log B)^C B^{-\epsilon'}$ by \eqref{eq:R5}.
The one potentially larger mean-curvature error is bounded by
\[
 C B^{1-m}\int_{\rho\asymp B}\mathcal F D
 \le C_R(1+\log B)
 B^{1-m}\left(\int_{\rho\asymp B}D^2\right)^{1/2}
             \left(\int_{\rho\asymp B}W_*^2\right)^{1/2}
 \le C_R(1+\log B)B^{-(1+\epsilon')/2}.
\]
We used $W_*\le C$, $\int D^2\le C B^{m-3}$, and
$\mathcal F\le C_R(1+\log B)W_*$.
Finally, choose $M_0$ sufficiently large and fixed. On
$d>M_0(L+\tau)$ we have $b(d)\ge c_1d$ and
$W_*\le Cd^{-\vartheta}e^{q_*c\tau-q_*c_1d}$.
Thus every fixed polynomial in $d,\rho$ and every
$e^{C\sqrt{\log(K(d)+1)}}$ occurring in the reach inequalities
has an integrable negligible tail there. This justifies truncating
those calculations to $d\le M_0(L+\tau)$. Here $M_0$ is chosen to
handle the finitely many polynomial degrees in these calculations;
it remains fixed through the successive limits $B\to\infty$ at fixed
$R$, then $R\to\infty$.
\subsection{The integrated reach inequality}

The preceding estimates pay for all cutoff derivatives and for the
large-reach tail. We now integrate the reach equation against the
radial weight, retaining the coefficients that will be compared with
stability.


For the rest of the reach calculation, unmarked derivatives are surface
derivatives. The definitions of $T_r$ and $\mathcal L$ are those above.
The hypersurface identity
$\Delta_\Sigma F=\tr_{T\Sigma}D^2F-H\nu\cdot\nabla F$
gives
\[
 \Delta\tau=\frac{m-2+2T_r-Hz\cdot\nu}{\rho^2},
 \qquad
 -\frac{\Delta\phi}{\phi}
 =\frac{(m-2)(mT_r-Hz\cdot\nu)}{\rho^2}.
\]

\begin{lemma}[Radial reach budget]\label{lem:radial-reach-budget}
With the fixed parameter choices above, for large $R$ and $B$,
\begin{equation}\label{eq:R6}
\begin{split}
 &\int_\Sigma\phi m_c^2W_*
 \left[
 \frac{3.44+24(\mathcal F/W_*)T_r}{\rho^2}
 +\frac{.012}{d}|\nabla d|^2
 +\frac{\lambda}{b'}|\nabla(c\tau-b(d))|^2
 \right]\dd S\\
 &\hspace{2em}\le .73\mathcal L+o_R(\mathcal A)+o_{B\to\infty;R}(1).
\end{split}
\end{equation}
\end{lemma}

\begin{proof}
Integrate \eqref{eq:R2} against $\phi m_c^2$ and move the surface
Laplacian onto this test. The preceding error estimates justify this
operation and control all derivatives of $m_H,m_B$.
The mean-curvature terms in the radial identities are at most
\[
 C\sup_{\rho\ge R}\frac{L+\tau}{\rho}
 \int\phi m_c^2W_*|H|=o_R(1)\mathcal L
\]
by \eqref{eq:R1}, \eqref{eq:radial-scale-separation}, and
$|H|\le .355I$. The term involving $J_1$ loses only
$o_R(1)\mathcal F_\tau/\rho^2$.

We spell out the source accounting in $W_*(\Delta d)_{\rm upper}$.
On a nonplateau branch where $ad\le a_0$,
\eqref{eq:C8} and $|\nabla d|\ge c_\Gamma>0$ show that
the sum of the $a^2d$ term and its gradient error is bounded by
$.008|\nabla d|^2/d$, on first choosing $a_0$ small and then the
uniform lower reach large. Where $ad>a_0$,
\eqref{eq:C7}--\eqref{eq:C8} instead bound the source by
\[
 C W_*a^4d^3e^{C\sqrt{\log(K(d)+1)}}.
\]
For $d\le M_0(L+\tau)\ll\rho$ its shell integral against $\phi$
is bounded by \eqref{eq:R4}, using $W_*\le d^{-\vartheta}$ and
\eqref{eq:M1}. Indeed, on $\rho\asymp s$, the reach power contributes
at most $C(L+\log s)^{3-\vartheta}$, the subpolynomial factor is
bounded using $K(d)\le K(Cs)$, and the product of $\phi$ with the
fourth-order shell bound has size $CK(s)s^{-2}$. The remaining tail
was treated above.
The gradient error on plateau branches is smaller than the same
$.008$ allowance.

The local mean curvature on a plateau costs $G_1|H(z)|$.
For its positive foot term use \eqref{eq:C10}.
On the truncated range the source and receiving radii have ratio
$1+o_R(1)$. Since $d(y)\le d(z)$ and
$|\partial_\tau\log W_*|\le cq_*$,
\[
 W_*(z)\le(1+o_R(1))W_*(y).
\]
If the receiving point lies in the rough collar, transfer first and
integrate over that receiving set. In fact \eqref{eq:C10} and the
uniformly diverging source reach imply $D(y)\to0$, so the foot is
good for the fixed threshold. The transfer Jacobian is bounded,
$W_*(z)\le d(z)^{-\vartheta}\le C$, $H_{N_y}(y)_+\le CD(y)\le C$,
and $\phi(z)\le(1+o_R(1))\phi(y)$. Positive incoming multiplicity
is at most one. Thus the charge is bounded by
\[
 C\int_{\substack{y\in\Sigma\setminus X\\
       \dist(y,Z)\le4\rho(y)^{\beta_0}}}\phi(y)\dd S_y.
\]
Its shell bound is $CK(s)s^{-2+2\alpha+n\beta_0}$, smaller than
\eqref{eq:R3}. No enlargement of the receiving collar by $d(z)$
is made: even when $d(z)$ is of order $L$, the area formula counts
the receiving set itself.
Otherwise the point is on the good surface. For
$a(y)d(z)>a_1$, use
$H_{N_y}(y)_+\le C/d(z)\le C_{a_1}a(y)^4d(z)^3$,
the uniformly bounded area distortion, and \eqref{eq:M1};
this is another error of type \eqref{eq:R4}. The reach power here
remains the source value $d(z)^3$: the truncated source and receiving
points lie in comparable shells, and the combined factor
$W_*(z)d(z)^3$ is at most $C(L+\log s)^{3-\vartheta}$ after transfer.
For $a(y)d(z)\le a_1$, choose $a_1$ and $\delta_2$ small so that
the area distortion is at most $1.01$.
There is exactly the incoming multiplicity bound of
\eqref{eq:C10}, namely one.

Spatial cutoff transfer uses, with fixed small $\epsilon_c>0$,
\[
 m_B(z)^2\le(1+\epsilon_c)m_B(y)^2+
 C_{\epsilon_c}(d/B)^2\mathbf1_{\{\rho\asymp B\}}.
\]
The latter term tends to zero by \eqref{eq:R5}, since the truncated
reach is at most $C(L+\log B)$ on this shell. The radial factors
$\phi$ have ratio $1+o_R(1)$ as well.
Consequently the sum of all local and incoming curvature charges is
at most
\[
 (1+o_R(1))G_1(.355)\bigl(1+1.01(1+\epsilon_c)\bigr)\mathcal L
 <.73\mathcal L .
\]
This leaves room for the radial mean-curvature terms already listed.

It remains to give the constant mass on the left. If $d\ge .8\tau$,
then $b(d)\le d$ and \eqref{eq:R1} gives
\[
 (4-o_R(1))\mathcal F_\tau/W_*
 \ge 4(.87)(.8/.87)^{.02}-o_R(1)>3.44 .
\]
The nonnegative $2T_r\mathcal F_\tau$ term may be discarded.
On the exceptional set $d<.8\tau$, $E^2\ge s^{-1.6\kappa}$ on
a shell $\rho\asymp s$. Thus the omitted weighted mass
$\int\phi W_*/\rho^2$ is at most $C K(s)s^{-4+1.6\kappa}$
by \eqref{eq:M1}. Its dyadic sum is negligible, since
$1.6\kappa<3$ and $K(s)\le C(s/R)K(R)$.
The gradient coefficient left from $\vartheta=.02$ is $.012$.
The second radial identity gives $m(m-2)=24$, proving \eqref{eq:R6}.
\end{proof}

\subsection{The local geometry for a global ambient test}

The remaining task is to recover almost all of $\mathcal L$ from ambient
stability. We will extend the amplitude $\psi$ normally and interpolate
between its values across a fixed-width middle band of each relevant
gap. This differs from the square-root gluing in
Section~\ref{sec:growth}: that construction proved the growth bound,
whereas we now use that bound to remove rough regions while interpolating
the sheet amplitudes themselves.

\begin{lemma}[Geometry of the radial interpolation regions]
\label{lem:radial-clean-geometry}
Fix $C_0<\infty$. For all sufficiently large $R$, the common graph and
projection conclusions of Lemma~\ref{lem:interface-log-windows} apply
within distance $C_0\log\rho(z)$ of every zero satisfying
$\dist(z,Z)>\rho(z)^{\beta_0}/3$. They apply in particular at every
base with $m_H\ne0$, uniformly for $\rho\ge R$ and independently of $B$.
The polynomial graph buffers accommodate every fixed logarithmic-power
lateral disk. Besides the exact identity $d(z)=G_1\ell(z)$ whenever a
distinct sheet is at logarithmic distance, where $\ell(z)$ is the
shortest adjacent Euclidean gap, the following conclusions hold.
\begin{enumerate}
 \item Two sheet distances within a fixed bounded amount of the least
 distance can compete only across an adjacent gap, on opposite sides
 of the point. Three sheets cannot so compete.
 \item Fix a base $z_i$ and one normal side. Its only possible partner
 in a middle band is the first adjacent graph on that side. Write
 $\ell_{ij}$ for the distance from the fixed base to that graph, and
 suppose $\ell_{ij}\le C_0\log\rho$. Before the ray reaches that graph,
 at logarithmic depths in the stated window, its two sheet distances
 $s_i,s_j$ satisfy
 \[
 s_i+s_j=\ell_{ij}+o_R(1),\qquad
 \frac{\dd}{\dd s_i}\frac{s_i-s_j}{2}=1+o_R(1).
 \]
 The ray traverses a fixed-width middle band at most once. The partner
 graph and $\ell_{ij}$ are fixed although the projected receiving
 point can move on that graph.
\end{enumerate}
All distance errors are additive and uniformly vanish. Buffered graph
families and their projections agree on overlaps.
\end{lemma}

\begin{proof}
Choose $0<\omega<\min(\alpha,\beta_0)$, put $N=\rho(z)$, and consider
$B_{2N^\omega}(z)$. Since $\rho$ is $1$-Lipschitz, its values there are
$N+O(N^\omega)$. For large $R$ this ball avoids $Z$, because
$2N^\omega<N^{\beta_0}/3$. Every zero in it therefore satisfies
\[
 D<\rho^{-\alpha}\le CN^{-\alpha}\le N^{-\omega}.
\]
Lemma~\ref{lem:interface-log-windows} applies on a fixed fractional
interior, decreasing $\omega$ if necessary to leave a margin. Every
fixed logarithmic-power disk and every fixed normal interval fit
inside that interior. Its errors are bounded by a positive inverse
power of $N$ times a fixed logarithmic power, so their supremum for
$N\ge R$ tends to zero. The sharp analytic estimates
\eqref{eq:L12}--\eqref{eq:L13} apply with these same premises after
$P$ and the depth constants are fixed.

We verify the extra conclusions used to count bands. In a frame over
one of the common tangent planes, unsigned sheet distances differ
additively by $o_R(1)$ from distances to the ordered graph intercepts.
Successive intercepts differ by at least $c_0\log\rho$. Among ordered
points on a line, two whose distances are within a fixed amount of
the least distance must bracket the point and be adjacent; a third
is farther by a diverging amount. The same assertions therefore hold
for the true sheets.

Now fix a base and its positive normal ray, putting the base at height
zero. The first adjacent graph on that side has intercept
$l=\ell_{ij}+o_R(1)$. Before reaching it, the graph estimates and
smooth distance gradients give
\[
 s_j=l-s_i+o_R(1),\qquad\frac{\dd s_j}{\dd s_i}=-1+o_R(1).
\]
These are the asserted sum and derivative formulas. The derivative of
$(s_i-s_j)/2$ is positive, hence a fixed-width band is traversed at most
once. A farther graph cannot be the partner before the first is reached,
because the first intervenes and is closer. After crossing that first
graph, it remains nearer than the base graph on the relevant logarithmic
ray. Thus the base cannot meet a second band on this side. The negative
side is identical. Finally, all descriptions concern the actual zero
pieces and their unique projections, so they agree on overlaps.
\end{proof}

\subsection{A single compactly supported test in physical space}

We now extend the sheet values $\psi$ to the ambient space. The interpolation is chosen to match the phase tail. In a gap of length $\ell$, \eqref{eq:L13} weights its normal derivative by a constant multiple of $e^{-\sqrt2\ell}\cosh^2(\sqrt2t)$. For a scalar interpolation $F$ from $1$ to $0$ on $[-A,A]$, Cauchy--Schwarz gives
\[
 1=\left|\int_{-A}^{A}F'\dd t\right|^2
 \le\left(\int_{-A}^{A}\cosh^2(\sqrt2t)|F'|^2\dd t\right)
       \left(\int_{-A}^{A}\sech^2(\sqrt2t)\dd t\right).
\]
Equality holds exactly when $F'$ is a negative constant times $\sech^2(\sqrt2t)$. Fixing the half-width at $A_{\rm sw}=3$, we therefore define
\begin{equation}\label{eq:radial-interpolant}
 \mathcal U(t)=
 \frac{\int_t^{3}\sech^2(\sqrt2u)\dd u}
      {\int_{-3}^{3}\sech^2(\sqrt2u)\dd u},\qquad -3\le t\le3.
\end{equation}
This coefficient is unrelated to $U=g^{-1}(v)$. Choose a large fixed depth constant $M$, and use buffered local graph families containing every candidate within distance $2M\log\rho+10$. More distant pieces cannot enter the bands below. For points satisfying
\[
 \dist(x,\Sigma)\le M\log\rho(x)+1,
 \qquad \dist(x,Z)>\tfrac12\rho(x)^{\beta_0},
\]
let $z_i(x)$ denote projection to the nearest sheet. Set
$\widetilde\eta(x)=\psi(z_i(x))$ unless its distance and the next smallest distance obey $|s_i-s_j|\le6$. In that band set
\begin{equation}\label{eq:radial-ambient-amplitude}
 \widetilde\eta(x)=\mathcal U(t)\psi(z_i(x))+
                  (1-\mathcal U(t))\psi(z_j(x)),
 \qquad t=(s_i-s_j)/2.
\end{equation}
Labels in this formula are local and may be chosen in either order.

All participating feet lie a logarithmic distance from $x$, so
$\dist(z_i,Z)>\rho(z_i)^{\beta_0}/3$ for large $R$. Thus
Lemma~\ref{lem:radial-clean-geometry} supplies the projection charts even when all projected amplitudes vanish. There is no triple competition.
Moreover $\mathcal U(-t)=1-\mathcal U(t)$, so exchanging the two labels leaves the formula unchanged. At $t=-3$ and $t=3$ it equals the appropriate single-sheet value. Near a zero it uses only the smooth projection, without differentiating that zero's unsigned distance. The overlap assertion of Lemma~\ref{lem:radial-clean-geometry} therefore makes $\widetilde\eta$ one locally Lipschitz function.

Let $\zeta$ be a Lipschitz function equal to one for
$\dist(x,\Sigma)\le M\log\rho(x)$ and zero for
$\dist(x,\Sigma)\ge M\log\rho(x)+1$, with bounded gradient. For example apply a fixed Lipschitz cutoff to
$\dist(x,\Sigma)-M\log\rho(x)$. Set $\eta=\zeta\widetilde\eta$ in the displayed region and zero elsewhere. Near its auxiliary boundary
$\dist(x,Z)=\rho(x)^{\beta_0}/2$, every participating foot lies at distance less than $\rho(z_i)^{\beta_0}$ from $Z$, hence has $m_H=0$. There is a full zero neighborhood before this boundary, so extension by zero remains locally Lipschitz.

If $\eta(x)\ne0$, some participating base has $|z_i|\le2B$, and
$|x-z_i|\le C(1+\log\rho(x))$. Consequently
\[
 |x|\le2B+C(1+\log(R+|x|)),
\]
which bounds $|x|$ for fixed $R,B$. A finite cover of this compact support by the local projection charts gives a finite Lipschitz constant. Uniformity of that constant as $B$ grows is unnecessary. We have constructed an admissible test for \eqref{eq:q-stability} without a global enumeration or a global bound on the number of sheets.



\subsection{The normal switching cost}

\begin{lemma}[Normal switching cost]\label{lem:radial-switch-cost}
With the preceding fixed parameter choices, for sufficiently large $R$
and uniformly in the subsequent outer radius $B$,
\begin{equation}\label{eq:R7}
 .895\mathcal L\le1.005\int_\Sigma|\nabla_\Sigma\psi|^2\dd S
                  +o_R(\mathcal A)+o_{B\to\infty;R}(1).
\end{equation}
The constructed ambient test has normalized switching cost at most
$.087853\mathcal L$, apart from negligible collar and cutoff terms.
\end{lemma}

\begin{proof}
We separate the terms whose coefficients tend to their planar values
from the normal derivative of the interpolation. This leaves a one-gap
calculation and the count that allows a positive charge only once per
base.

\paragraph{Defect mass, tangential derivatives, and cutoffs.}

Over supported bases the fixed tubes $|s|\le P$ are jointly injective for large $R$: any intersecting segments have bases in one buffered chart, where separation excludes different sheets and uniqueness of projection excludes two bases on one sheet. They have Jacobian $1+o_R(1)$. The base is nearest there, $\zeta=1$, and
$\eta(z+s\nu(z))=\psi(z)$. Hence
\begin{equation}\label{eq:radial-test-mass}
 \int q^3D\eta^2\dd x\ge(1-o_R(1))\sigma\mathcal L.
\end{equation}

In a single-projection region the gradient of $\psi(z_i(x))$ is
$(\mathrm{Id}+sA_i)^{-1}\nabla_\Sigma\psi$. Its metric and volume factors are $1+o_R(1)$ at every relevant logarithmic depth. Where $i$ is nearest, the distance to $\Sigma$ equals $|s|$. Fixed-depth profile convergence and then the uniform phase tail in \eqref{eq:phase-decay} give
\[
 \int_{\text{nearest part of the ray}}q^2\dd s\le\sigma+o_R(1).
\]
For precision, first choose a large fixed depth to make the exponential tail arbitrarily small, then increase $R$ to use profile convergence on the fixed interval. This proves a uniform limsup without assuming convergence on an expanding interval.

In a band, differentiation of the two base amplitudes gives
$\mathcal U\nabla(\psi\circ z_i)+(1-\mathcal U)\nabla(\psi\circ z_j)$. Convexity bounds its square by the corresponding convex combination of gradient squares. Both sheet distances are at least a fixed multiple of $\log\rho$, so the $q^2$ mass of a fixed-width band tends uniformly to zero. Normal parametrization from each base, using the single traversal in Lemma~\ref{lem:radial-clean-geometry}, bounds the integral of these tangential terms by
$o_R(1)\int|\nabla_\Sigma\psi|^2$. Young's inequality with a fixed small parameter on the switch derivative changes its square by a factor $1+\epsilon_s$; its reciprocal factor on this tangential part is harmless. Combining all such terms gives
\begin{equation}\label{eq:radial-tangent-cost}
 (1+o_R(1))\sigma\int_\Sigma|\nabla_\Sigma\psi|^2\dd S.
\end{equation}

On the support of $\nabla\zeta$ the phase bound gives
$q^2\le C\rho^{-2M(\sqrt2-\gamma)}$. Also
$\widetilde\eta^2\le C\rho^{-4}$, because $W_*\le d^{-\vartheta}$ is uniformly bounded on supported bases and their radii differ by $O(\log\rho)$. Ambient shell volume is at most $C\rho^7$. Thus the depth cost is bounded by
\[
 C\sum_{j\ge0}s_j^{3-2M(\sqrt2-\gamma)}=o_R(\mathcal A)
\]
for sufficiently large fixed $M$. Its cross term uses another fixed Young factor. The only cost still needing a fixed fraction of $\mathcal L$ is
\begin{equation}\label{eq:radial-switch-integral}
 S_{\rm sw}=\int_{\rm bands}q^2\mathcal U'(t)^2
                  |\psi(z_i)-\psi(z_j)|^2\dd x,
\end{equation}
since $|\nabla t|\le1$.

If either projected $m_H$ is not one, a nonzero projected base is in a fixed logarithmic enlargement of the $4\rho^{\beta_0}$ collar of $Z$. Parametrize by the normal rays of the nonzero bases. Each base has at most two bands, with bounded widths and polynomially bounded logarithmic Jacobian factors. The integrand is at most $C\rho^{-4}$ times a fixed logarithmic power. The collar bound in Lemma~\ref{lem:radial-errors} and the summation \eqref{eq:R3} make this portion negligible. This argument counts rays, not globally named sheets.

On the remaining bands $m_H=1$ at both feet. Assign each physical point to a foot $i$ for which the uncut amplitude $V_i$ is larger, breaking equality measurably. This partitions the nonnegative integral \eqref{eq:radial-switch-integral}; it does not redefine $\eta$, so no derivative of the assignment occurs. With $r=V_j/V_i\in(0,1]$,
\begin{equation}\label{eq:radial-cutoff-split}
 |m_B(z_i)V_i-m_B(z_j)V_j|^2
 \le(1+\epsilon_s)m_B(z_i)^2V_i^2(1-r)^2
      +C_{\epsilon_s}V_i^2|m_B(z_i)-m_B(z_j)|^2.
\end{equation}
The last term is supported on $\rho\asymp B$ when $B$ is large at fixed $R$, and its cutoff difference is at most $C\log B/B$. Normal parametrization and \eqref{eq:R5} therefore bound its integral by
$C_R(1+\log B)^kB^{-\epsilon'}\to0$.

\paragraph{One fixed base and one adjacent side.}

Fix an assigned base $z_i$ and one normal side that contributes a band. Its partner is the first adjacent graph on that side, fixed throughout the band. Let $\ell_{ij}$ be the shortest Euclidean distance from this fixed base to that graph, and put $\ell_i=d(z_i)/G_1$. By Lemma~\ref{lem:radial-clean-geometry}, $\ell_i$ is the shortest adjacent gap, so
$u:=\ell_{ij}-\ell_i\ge0$. The projected foot $z_j$ moves along its graph as the normal depth changes; $\ell_{ij}$ does not.

We justify the additive error for the moving foot. Put $z_i=0$ and
its normal ray on the vertical axis, so $x=(0,s_i)$ and $s_i$ is the
actual normal depth. Write the partner graph as $x_7=f(y)$, with
common-plane slope bounded by $\theta_N=o(1)$ at a polynomial rate,
where $N=\rho(z_i)$. The stationary equation for projection gives
$|y_j|\le C\theta_N\log N$ and
$f(y_j)=f(0)+O(\theta_N^2\log N)$. The logarithmic lower separation
then gives
\[
 |z_j-z_i|=f(0)+O(\theta_N^2\log N)
          =\ell_{ij}+o_R(1),
\]
after choosing the positive vertical side; the negative side is
identical with absolute heights. The second equality also follows by
applying the same projection equation to the base $z_i$ itself.
The chord makes an unoriented angle tending to zero with both graph
normals, so its
penalty at $z_j$ is on the plateau. It is an admissible receiving-reach
competitor, giving
\begin{equation}\label{eq:radial-receiving}
 d(z_j)\le G_1(\ell_{ij}+o_R(1)),\qquad d(z_i)=G_1\ell_i.
\end{equation}
The radial factors satisfy $h_j/h_i=1+o_R(1)$ and
$|\tau_j-\tau_i|\le C\log\rho/\rho$. By
\eqref{eq:radial-weight-slopes}, monotonicity in $d$, and the uniform divergence of the supported reaches, there is $\epsilon_R\to0$ such that, throughout all assigned bands,
\begin{equation}\label{eq:radial-amplitude-ratio}
 \frac{V_j}{V_i}\ge(1-\epsilon_R)
     \exp[-(q_*G_1/2+\epsilon_R)u].
\end{equation}
Here both feet have $m_H=1$ and thus diverging reach. To verify the inequality when $d(z_j)>d(z_i)$, integrate the logarithmic derivative in $d$ from \eqref{eq:radial-weight-slopes} from $d(z_i)$ up to the upper bound in \eqref{eq:radial-receiving}, and use its $\tau$-derivative bound to change $\tau$. If $d(z_j)\le d(z_i)$, decreasing reach only increases $W_*$, so the same lower bound holds. The additive distance error in \eqref{eq:radial-receiving} enters a multiplicative factor tending to one. This proves \eqref{eq:radial-amplitude-ratio} uniformly for $\rho\ge R$.

The sharp phase input and Lemma~\ref{lem:radial-clean-geometry} give
\begin{equation}\label{eq:radial-q-band}
 q^2\le64\left(1+\frac{.001}{4}\right)^2(1+o_R(1))
       e^{-\sqrt2\ell_{ij}}\cosh^2(\sqrt2t).
\end{equation}
Indeed $e^{-\sqrt2s_i}+e^{-\sqrt2s_j}
=2e^{-\sqrt2(s_i+s_j)/2}\cosh(\sqrt2t)$, and its sum of distances differs additively from $\ell_{ij}$ by $o_R(1)$.
The finite-width integral is exact:
\begin{equation}\label{eq:radial-finite-integral}
 \mathcal U'(t)=-\frac{\sech^2(\sqrt2t)}{\sqrt2\tanh(3\sqrt2)},
 \qquad
 \int_{-3}^{3}64\cosh^2(\sqrt2t)\mathcal U'(t)^2\dd t
 =\frac{32\sqrt2}{\tanh(3\sqrt2)}.
\end{equation}
The change from normal depth to $t$ and the normal volume Jacobian are $1+o_R(1)$. Therefore, after division by the defect mass at this base,
\[
 \sigma I_P(z_i)m_B(z_i)^2V_i^2,
\]
the contribution of this side is bounded, up to the displayed finite profile factors and vanishing errors, by
\begin{equation}\label{eq:radial-one-edge}
 \frac{1}{\tanh(3\sqrt2)}\frac{\sqrt2}{\sigma}
 e^{-\sqrt2u}(1-e^{-q_*G_1u/2})^2.
\end{equation}
We used $I_P(z_i)\ge31.999e^{-\sqrt2\ell_i}$; replacing $31.999$ by $32$ in \eqref{eq:radial-one-edge} contributes the explicit factor $32/31.999$ below. A base with $m_B=0$ contributes no principal term in \eqref{eq:radial-cutoff-split}, so division by zero is never needed.

The vanishing error in this normalized assertion is uniform even if $u$ grows with $\log\rho$. For any fixed $U_0$, the lower bound in \eqref{eq:radial-amplitude-ratio} converges uniformly on $0\le u\le U_0$. On $u>U_0$ both the actual loss and its proposed bound are bounded by a fixed constant times $e^{-\sqrt2u}$, because $0\le1-r\le1$. First make this tail small by choosing $U_0$, then increase $R$. In particular the side with $u=0$ has loss tending to zero: \eqref{eq:radial-amplitude-ratio} then gives $1-r\le\epsilon_R$.

The count of sides now matters. There is at most one first adjacent graph on each side, and the normal ray crosses its band at most once. If both sides contribute, at least one realizes the nearest adjacent distance $\ell_i$; if only one side contributes and the nearest side is different, that single contributing side is the only possible positive charge. Thus at most one side needs the positive supremum of \eqref{eq:radial-one-edge}. Restricting to the measurable part assigned to $z_i$ can only decrease the integral. No factor of two multiplies this supremum. Figure~\ref{fig:radial-interpolation} summarizes this count.

\begin{figure}[ht]
\centering
\begin{tikzpicture}[x=1cm,y=.75cm,>=Latex,font=\small]
\fill[blue!10](-3,-.95) rectangle (3,-.65);
\fill[orange!16](-3,1.1) rectangle (3,1.4);
\draw[thick](-3,-1.6).. controls(-1,-1.62)and(1,-1.58)..(3,-1.55);
\draw[thick](-3,0)..controls(-1,-.03)and(1,.03)..(3,.06);
\draw[thick](-3,2.5)..controls(-1,2.48)and(1,2.53)..(3,2.57);
\draw[dashed,->](0,-1.9)--(0,2.9) node[above]{normal line through $z_i$};
\fill(0,0)circle(2pt)node[below right]{$z_i$};
\draw[<->](-2.4,-1.6)--(-2.4,0)node[midway,left]{$\ell_i$};
\draw[<->](2.4,.06)--(2.4,2.55)node[midway,right]{$\ell_{ij}$};
\node[anchor=west]at(3.5,-.8){nearest band};
\node[anchor=west]at(3.5,1.25){possible charged band};
\end{tikzpicture}
\caption{A schematic local cross-section, not to scale. The colored strips represent fixed-width middle bands. On the part of a band assigned to $z_i$ as the larger uncut amplitude, the shorter gap $\ell_i=d(z_i)/G_1$ has vanishing normalized loss. Here $\ell_{ij}$ is the distance from $z_i$ to its first adjacent sheet on the other side. Only that side can require a positive fixed allowance. The proof uses true normal projections and does not assume globally parallel sheets.}
\label{fig:radial-interpolation}
\end{figure}

\paragraph{The finite tolerance budget.}

Since $\sqrt2/\sigma=3/4$, define
\begin{equation}\label{eq:radial-loss-number}
 C(G_1)=\frac34\sup_{u\ge0}
       e^{-\sqrt2u}(1-e^{-1.1G_1u})^2.
\end{equation}
Set $x=e^{-1.1G_1u}$ and $a_*=\sqrt2/(1.1G_1)$. Differentiation of
$x^{a_*}(1-x)^2$ shows its only interior maximum is at
$x=a_* /(a_*+2)$, while the endpoint values are zero. Hence
\begin{equation}\label{eq:radial-maximum}
 C(G_1)=\frac34
 \left(\frac{a_*}{a_*+2}\right)^{a_*}
 \left(\frac2{a_*+2}\right)^2.
\end{equation}
The defining expression is nondecreasing in $G_1$, so the largest permitted value occurs at $1.001$, and direct evaluation gives
$C(G_1)<.083272621$.

Use Young factors at most $1.001$ for the band derivative, depth cutoff, and cutoff difference. After those fixed choices, increase $R$ until the normal Jacobian, the inverse derivative $\dd s_i/\dd t$, and the exponential of the additive distance error are each at most $1.001$. After all these finite factors have been fixed, the uniform estimate following \eqref{eq:radial-one-edge} allows a normalized error of at most $.002$ on each side, including the nearest side, \emph{after multiplication by those factors}. This is legitimate because the unamplified uniform error can be made arbitrarily small.  The principal normalized switching contribution is consequently at most
\begin{equation}\label{eq:radial-finite-budget}
 \frac{.083272621}{\tanh(3\sqrt2)}
 \left(1+\frac{.001}{4}\right)^2
 \frac{32}{31.999}(1.001)^6+.004
 <.087853<.100.
\end{equation}
Every factor is finite: the band width was fixed at $3$, and all analytic tolerances were fixed before $R$ increased. In particular the proof does not use an infinite-width switch and silently discard its tails.

Choose the depth constant $M$ large enough for its negligible cost before the final increase of $R$. Then \eqref{eq:radial-test-mass} is at least
$.998\sigma\mathcal L$, and the sheet-gradient costs in
\eqref{eq:radial-tangent-cost}, including the fixed Young factor for the depth cutoff, are at most
$1.005\sigma\int|\nabla_\Sigma\psi|^2$.
The other Young factors multiply only the switch term or a tangential band term tending to zero. The rough collar and outer cutoff errors have already been shown negligible. Substitution into \eqref{eq:q-stability} and division by $\sigma$ give
\[
 (.998-.100)\mathcal L
 \le1.005\int_\Sigma|\nabla_\Sigma\psi|^2\dd S
       +o_R(\mathcal A)+o_{B\to\infty;R}(1).
\]
As $.998-.100>.895$, this proves \eqref{eq:R7}. All estimates were made at finite $B$ with coefficients uniform in $B$; the outer limit therefore precedes the limit $R\to\infty$ exactly as required by the radial argument.
\end{proof}

\subsection{The coefficient comparison and the contradiction}

The reach and stability inequalities now concern the same sheet
amplitude. Removing the derivative of its radial factor exposes the
coefficients directly; their strict differences will leave a positive
multiple of $\mathcal F T_r/\rho^2$.

The cross term against $h$ can be removed exactly:
\[
 \int_\Sigma|\nabla(hm_c\sqrt{W_*})|^2
 =\int_\Sigma\phi|\nabla(m_c\sqrt{W_*})|^2
       -\int_\Sigma h\Delta h\,m_c^2W_* ,
 \qquad
 -\frac{\Delta h}{h}
 =\frac{4+8T_r-2Hz\cdot\nu}{\rho^2}.
\]
All terms have compact support in the regular surface.
The $H$ term is $o_R(1)\mathcal L$ by \eqref{eq:R1} and
\eqref{eq:radial-scale-separation}. The derivatives of $m_c$ have
negligible cost by Lemma~\ref{lem:radial-errors}; Young's inequality
costs an arbitrarily small fixed relative loss in the other
gradient term. Since
\[
 \frac{.73(1.005)}{.895}<.820,\qquad
 \nabla\log W_*=
 \lambda\nabla(c\tau-b(d))-\frac{\vartheta}{d}\nabla d,
\]
\eqref{eq:R7}, after absorbing $o_R(1)\mathcal L$, bounds
$.73\mathcal L$, up to negligible errors, by
\[
 .835\int_\Sigma\phi m_c^2W_*
 \left[
 \frac{4+8T_r}{\rho^2}
 +\frac{1.03}{4}
 \left|\lambda\nabla(c\tau-b(d))
                 -\frac{\vartheta}{d}\nabla d\right|^2
 \right]\dd S .
\]
Every comparison with \eqref{eq:R6} is strict. For the gradient
square use $|a-b|^2\le1.5|a|^2+3|b|^2$ and
\[
 .835\frac{1.03}{4}\,1.5\lambda^2
 \le .70955\lambda<.72\,\frac{\lambda}{b'}
 \quad(0<\lambda\le2.2,\ 0<b'\le1).
\]
Its other term is
$.835(1.03)3\vartheta^2/(4d^2)$, smaller than $.012/d$
uniformly on the supported set for large $R$.
Where $\lambda=0$ the first comparison is simply unnecessary.
The constant radial terms satisfy $.835\cdot4=3.34<3.44$.
For the $T_r$ term use $\mathcal F/W_*\ge1/(2.2+o_R(1))$:
\[
 24\frac{\mathcal F}{W_*}-.835\cdot8
 \ge\bigl(24-.835\cdot8(2.2+o_R(1))\bigr)
                 \frac{\mathcal F}{W_*}
 \ge5\frac{\mathcal F}{W_*}.
\]
Subtracting and discarding the other nonnegative remainders yields
\begin{equation}\label{eq:R8}
 \int_\Sigma\phi m_c^2\mathcal F\,\frac{T_r}{\rho^2}\dd S
 \le o_R(\mathcal A)+o_{B\to\infty;R}(1).
\end{equation}
Every absorption of $o_R(1)\mathcal L$ above is performed at finite
$B$, with its coefficient uniform in $B$. Thus the exhaustion does
not presume that $\mathcal L$ has a finite limit as $B\to\infty$.

It remains to find enough surface area where the weight has not
decayed. The chosen center supplies energy of order $K(R)R^m$.
We first transfer that energy to regular zero-surface area, then use
injective normal tubes to exclude a substantial portion at large reach.

\begin{lemma}[Area at small reach]\label{lem:radial-area}
For all sufficiently large $R$, the set
\[
 \{z\in\Sigma\cap B_{2R}:m_H(z)=1,\ d(z)\le L/2\}
\]
has area at least $c K(R)R^m$, with $c>0$ independent of $R$ and
of the selected center.
\end{lemma}

\begin{proof}
The energy in $B_R$ is at least $c K(R)R^m$ by the center choice.
We first transfer a fixed fraction of it to good surface area.
Let $\mathcal Z$ be a fixed-scale separated net on all of $\Sigma$.
Uniform gradient bounds give a fixed positive lower energy in a
fixed small ball at every zero, including singular zeros.
Packing, the definition of $K$, and ball coverings show that the
number of net points in $B_{R+j+2}$ is at most
\[
 C K(R)R^m(1+j/R)^{m+1}.
\]
By \eqref{eq:phase-decay} and Modica's inequality, the energy at
points of $B_R$ whose distance from $\Sigma$ is greater than $P_0$
is therefore at most
\[
 C K(R)R^m\sum_{j\ge P_0}
 (1+j/R)^{m+1}(j+2)^n e^{-c_1j}.
\]
Choose $P_0$ large and fixed so this is a small fraction of the
initial energy. Energy within the enlarged rough collar is also
a negligible fraction: its density is bounded and its volume in
$B_{2R}$ is $O(K(R)R^{m-4+2\alpha+n\beta_0})$.
Enlarge the collar by $2P_0$ without changing this estimate.
Every remaining point lies within $P_0$ of a zero $z$ in $B_{2R}$
with $m_H(z)=1$. Such bases have $D\to0$ uniformly. Their fixed
normal tubes are regular with bounded Jacobian, and cover the
points in question by nearest projection.
Their energy is at most a fixed constant times their base area
by the area formula and the bounded energy density. Thus
\[
 \cH^m(\Sigma\cap B_{2R}\cap\{m_H=1\})\ge c K(R)R^m.
\]


The large-reach part can now be removed with the jointly injective tube
estimate from Remark~\ref{rem:reach-ordinary-tube}, applied with threshold
$L/2$. More explicitly, its normal map is injective simultaneously over
all bases with $d>L/2$ and on both normal sides for $|t|<c_2L$,
where $c_2>0$ depends only on the fixed penalty. Its Jacobian has a
fixed positive lower bound. Since $L\le R$, the images of bases in
$B_{2R}$ lie in a ball of radius at most $3R$. Consequently
\[
 \cH^m(\Sigma\cap B_{2R}\cap\{d>L/2\})
 \le C R^{m+1}/L
 =C K(R)R^m/\sqrt{K(R)}=o(K(R)R^m).
\]
Subtracting this from the preceding good-area lower bound proves the
lemma.
\end{proof}

Take $B>3R$ in \eqref{eq:R8}. On the set in
Lemma~\ref{lem:radial-area}, $m_c=1$, $\rho\le\sqrt5R$ and
$T_r\ge1/5$. Since $d\le L/2$ and $f_*(c\tau)=1$,
\[
 \mathcal F(\tau,d)\ge\int_{L/2}^L t^{-\vartheta}\dd t
 =\frac{1-2^{-(1-\vartheta)}}{1-\vartheta}L^{1-\vartheta}.
\]
Also $\phi T_r/\rho^2\ge cR^{-m}$ on $B_{2R}$.
The left side of \eqref{eq:R8} is consequently at least
$cK(R)L^{1-\vartheta}=c\mathcal A$.
First let $B\to\infty$ at fixed $R$ and then let $R\to\infty$.
This contradicts \eqref{eq:R8}. The unbounded-density alternative
is impossible.

\begin{proof}[Completion of Theorem~\ref{thm:stable}]
If a solution attains either endpoint of $[-1,1]$, it is the corresponding
constant by the strong maximum principle. The remaining constant,
zero, is unstable. Indeed, for any nonzero $\chi\in C_c^1(\R^7)$,
the test $\chi_R(x)=\chi(x/R)$ gives
\[
 Q_0(\chi_R)=R^5\int|\nabla\chi|^2-R^7\int\chi^2<0
\]
for large $R$.
For a nonconstant stable solution with values in $(-1,1)$,
if $D$ vanishes anywhere, the equality case in
Section~\ref{sec:diffuse} gives a planar heteroclinic.
Otherwise $D>0$ everywhere. If its centered density is bounded,
Theorem~\ref{thm:finite} applies. If that density were unbounded,
then $K(R)\to\infty$ and Sections~\ref{sec:reach}--\ref{sec:growth},
followed by the contradiction just proved, would apply.
Thus this last case cannot occur, and Theorem~\ref{thm:finite}
gives precisely the claimed constants and planar profiles.
\end{proof}
